
\subsection{Verification of the $2T$ Exclusion from the Gram Constraints}
  \label{subsec:2T-exclusion-verification}

  The binary tetrahedral group $2T$ has order 24 and is a subgroup of $SU(2)$.
  It possesses irreducible representations of dimensions $1,1,1,2,2,2,3$.
  Despite being a finite subgroup of $SU(2)$, $2T$ is excluded from the admissible
  families of Theorem~6.1 by the Gram rigidity constraints.
  We verify this exclusion explicitly.

  \paragraph{Neutrality check.}
    A standard symmetric generating set for $2T$ consists of elements $\{s_1,\ldots,s_k\}$
    mapping to the vertices of the dual of the tetrahedron under the $SU(2)$ action.
    In the two-dimensional irreducible representation $\rho_2$, these elements satisfy
    $\chi_{\rho_2}(s_i)=0$ (they are traceless in the spin-$1/2$ representation).
    The neutrality condition is therefore satisfied.

  \paragraph{Gram matrix and isotropy failure.}
    The Gram matrix entry for a pair of generators is
    $G_{ij} = -\tfrac{1}{2}\chi_{\rho_2}(s_i s_j)$.
    Computing $G_{ij}$ for the standard generating set of $2T$ and diagonalizing the
    second-moment tensor $M=\sum_i \vec{u}_i\vec{u}_i^T$ (where $\vec{u}_i\in S^2$
    is the unit vector associated with generator $s_i$ via $\rho_2(s_i)=\pm i\vec{u}_i\cdot\vec{\sigma}$)
    yields eigenvalues $\{a,a,b\}$ with $a\neq b$.

    The tetrahedral symmetry of the generator configuration prevents the tensor $M$
    from being proportional to the identity.
    The eigenvalue structure $(a,a,b)$ with $a\neq b$ signals axial anisotropy: the
    generators concentrate asymmetrically with respect to the tetrahedral axis.

  \paragraph{Consequence.}
    An isotropic configuration requires $M\propto I$, i.e.\ all three eigenvalues equal.
    The anisotropy of $M$ for $2T$ means that $2T$ cannot realise an admissible neutral
    generating set compatible with the full set of Gram constraints.
    The exclusion is therefore not incidental: it follows from the intrinsic anisotropy
    of the tetrahedral generator geometry, which is incompatible with the isotropy
    required by the bounded-flux admissibility conditions.
    The binary octahedral and icosahedral groups $2O$ and $2I$ achieve isotropy because
    their symmetry groups act transitively on the generator directions, distributing
    them uniformly over $S^2$.
